\chapter{Аналитический раздел}

\section{Алгоритм AES}

AES (англ. \textit{Advanced Encryption Standard}; также \textit{Rijndael})~---~симметричный алгоритм блочного шифрования (размер блока 128 бит, ключ 128/192/256 бит)~\cite{aes}.

Далее опишем основные этапы данного алгоритма.

\subsection{Подготовка ключа}
Начальные ключи расширяются для создания ключей раундов. 
Количество раундов зависит от длины ключа: 10 раундов для 128 бит, 12 для 192 бит и 14 для 256 бит.

\subsection{Начальный раунд}
\texttt{AddRoundKey:} каждый байт блока данных складывается по модулю 2 с соответствующим байтом раундового ключа.

\subsection{Основные раунды}
Для каждого раунда, кроме последнего~\cite{aes}:
\begin{enumerate}
	\item \texttt{SubBytes:} каждый байт блока заменяется на соответствующий байт из S-буквенной таблицы (замена);
	\item \texttt{ShiftRows:} ряды блоков сдвигаются влево на разное количество байт;
	\item \texttt{MixColumns:} каждая колонка блока обрабатывается, чтобы смешать данные;
	\item \texttt{AddRoundKey:} выполняется XOR с новым раундовым ключом.
\end{enumerate}

\section{PCBC}

Эрсам, Мейер, Смит и Такман изобрели режим работы цепочки блоков шифрования (CBC) в 1976 году.
В режиме CBC каждый блок открытого текста перед шифрованием связывается с предыдущим блоком зашифрованного текста.
Таким образом, каждый блок зашифрованного текста зависит от всех блоков открытого текста, обработанных до этого момента.
Чтобы сделать каждое сообщение уникальным, в первом блоке должен использоваться вектор инициализации.

Если первый блок имеет индекс 1, математическая формула для шифрования CBC выглядит так:
\begin{equation}
	C_i = E_k(P_i \oplus C_{i-1})),
	C_0 = IV
\end{equation}
Для дешифрования:
\begin{equation}
	P_i = D_k(C_i) \oplus C_{i-1},
	C_0 = IV
\end{equation}

CBC является наиболее часто используемым режимом работы.
Его основные недостатки заключаются в том, что шифрование является последовательным и что сообщение должно быть дополнено до размера, кратного размеру блока шифрования.
Один из способов решения этой последней проблемы --- метод, известный как кража зашифрованного текста. 
При расшифровке с неправильным IV первый блок открытого текста будет повреждён, но последующие блоки открытого текста будут корректными.
Это связано с тем, что каждый блок шифруется по схеме XOR с зашифрованным текстом предыдущего блока, а не с открытым текстом, поэтому нет необходимости расшифровывать предыдущий блок перед использованием его в качестве IV для расшифровки текущего блока.
Это означает, что блок открытого текста можно восстановить по двум соседним блокам зашифрованного текста.
Следовательно, расшифровку можно выполнять параллельно.

PCBC был разработан таким образом, чтобы при расшифровке, как и при шифровании, небольшие изменения в зашифрованном тексте распространялись на весь текст.
В режиме PCBC каждый блок открытого текста перед шифрованием складывается операцией XOR с предыдущим блоком открытого текста и предыдущим блоком зашифрованного текста.
Как и в режиме CBC, в первом блоке используется вектор инициализации.
В отличие от CBC, при расшифровке PCBC с неправильным IV (вектором инициализации) все блоки открытого текста будут повреждены~\cite{pcbc}.
